\begin{proof}[Proof of \cref{obj:thm:unrestricted-fixed-q-minimax-frontier}]
Fix \(0<q<1/2\). For \(n,d\ge1\), put
\[
\mathfrak h_{n,d}^{\circ}
:=
\frac1n+\left(\frac{d}{n\log(e+n)}\right)^2,
\qquad
\mathfrak h_{n,d}:=\min\{1,\mathfrak h_{n,d}^{\circ}\}.
\]
For a bounded estimator kernel \(\mathsf T\) and a full-data law \(P\), write
\[
\mathcal R_P(\mathsf T)
:=
\int\!\int_{[-1,1]}
(h-\tau(P))^2\,\mathsf T(o,dh)\,
\mathcal L_P(\mathbf O_n)(do).
\]
Thus \(\mathcal R_P(\mathsf T)\) integrates both the sample and the estimator draw.

\begin{enumerate}
\item
The deterministic estimator in \cref{obj:def:mixed-count-estimator} defines an admissible decision kernel
\[
\mathsf T_{\mathrm{mix}}(o,dh)
:=
\delta_{\widehat\tau^{\mathrm{mix}}_{n,d,q}(o)}(dh).
\]
By \cref{obj:thm:unrestricted-mixed-count-upper}, there is a universal constant \(\overline C_{\mathrm{mix}}>0\) such that
\[
0\le \mathcal R_P(\mathsf T_{\mathrm{mix}})
\le \overline C_{\mathrm{mix}}\,r_{n,d,q}
\qquad
\text{for every }P\in\mathcal M^+_{n,d,q}.
\]
Taking the supremum over the model and then using this kernel in the defining infimum gives
\[
\mathfrak R^+_{n,d,q}
\le
\sup_{P\in\mathcal M^+_{n,d,q}}
\mathcal R_P(\mathsf T_{\mathrm{mix}})
\le
\overline C_{\mathrm{mix}}\,r_{n,d,q}.
\]
If the model class is empty, its real-valued worst-case risk is taken to be zero, and the same inequality follows from
\[
r_{n,d,q}
=
\min\left\{1,\frac1{nq}
+\left(\frac{d}{nq\log(e+nq)}\right)^2\right\}
\ge0.
\]
% lean: CausalSmith.Stat.MarRareqLogfrontier.unrestricted_fixed_q_frontier_upper

\item
We next convert this upper bound to the fixed-\(q\) rate. Define
\[
\beta_q:=\frac{1-\log q}{q},
\qquad
N:=nq,
\qquad
\ell:=\log(e+N),
\qquad
\ell_{n,1}:=\log(e+n).
\]
Throughout the present domain,
\[
N>0,\qquad \ell\ge1,\qquad \ell_{n,1}>0,
\qquad \log q\le0,\qquad \beta_q\ge1.
\]
Since
\[
e+n\le\frac{e+N}{q},
\]
monotonicity of the logarithm yields
\[
\ell_{n,1}\le\ell-\log q
\le(1-\log q)\ell.
\]
The second inequality follows from
\[
(-\log q)(\ell-1)\ge0.
\]
Consequently,
\[
n\ell_{n,1}
\le n(1-\log q)\ell
=\beta_q N\ell.
\]
Moreover,
\[
q\beta_q=1-\log q\ge1,
\qquad
q\beta_q^2=(1-\log q)\beta_q\ge1,
\]
so
\[
\frac1N\le\frac{\beta_q^2}{n}.
\]
The positive denominator comparison also gives
\[
0\le\frac{d}{N\ell}
\le\beta_q\frac{d}{n\ell_{n,1}},
\qquad
\left(\frac{d}{N\ell}\right)^2
\le\beta_q^2
\left(\frac{d}{n\ell_{n,1}}\right)^2.
\]
Adding these bounds,
\[
\frac1N+\left(\frac{d}{N\ell}\right)^2
\le\beta_q^2\mathfrak h_{n,d}^{\circ}.
\]
Truncation therefore gives
\[
r_{n,d,q}
\le\min\{1,\beta_q^2\mathfrak h_{n,d}^{\circ}\}
\le\beta_q^2\mathfrak h_{n,d}.
\]
Indeed, if \(\mathfrak h_{n,d}^{\circ}\le1\), the last inequality is
\[
\min\{1,\beta_q^2\mathfrak h_{n,d}^{\circ}\}
\le\beta_q^2\mathfrak h_{n,d}^{\circ}
=\beta_q^2\mathfrak h_{n,d};
\]
if \(\mathfrak h_{n,d}^{\circ}\ge1\), it is
\[
\min\{1,\beta_q^2\mathfrak h_{n,d}^{\circ}\}
\le1\le\beta_q^2
=\beta_q^2\mathfrak h_{n,d}.
\]
Define the positive fixed-\(q\) constant
\[
\overline C_{\mathrm{poly}}(q)
:=\overline C_{\mathrm{mix}}\beta_q^2.
\]
The upper bound from the preceding step now implies
\[
\mathfrak R^+_{n,d,q}
\le\overline C_{\mathrm{mix}}r_{n,d,q}
\le\overline C_{\mathrm{poly}}(q)\mathfrak h_{n,d}.
\]
% lean: CausalSmith.Stat.MarRareqLogfrontier.frontierRate_le_fixedQRate_mul; CausalSmith.Stat.MarRareqLogfrontier.unrestricted_fixed_q_polynomial_upper

\item
For the lower bound, invoke the published fixed-positivity lower bound on the control-zero observational subclass
\citep[Theorem~2, Section~3.2, p.~9, with its proof in Section~C.5, pp.~39--44]{ZengBalakrishnanHanKennedy2026}.
Its truncated, all-\(n,d\) comparison is supplied by
\cref{obj:thm:zeng-fixed-positivity-polynomial-frontier}: at \(\epsilon=q\), there is a constant \(\underline c_Z(q)>0\) such that
\[
\underline c_Z(q)\mathfrak h_{n,d}
\le\mathfrak R^Z_{n,d,q}
\qquad(n,d\ge1).
\]

We transfer this lower bound to the unrestricted experiment. For an observational estimator kernel \(\widehat\theta\), define
\[
\mathcal R^Z_{P_Z}(\widehat\theta)
:=
\int\!\int_{[-1,1]}
(h-\theta_Z(P_Z))^2\,
\widehat\theta(\mathbf z,dh)\,
\mathcal L_{P_Z}(\mathbf Z_n)(d\mathbf z).
\]
Given any admissible unrestricted kernel \(\mathsf T\), define its observational pullback by
\[
\widehat\theta_{\mathsf T}(\mathbf z,\mathcal V)
:=
2^{-n}
\sum_{\mathbf u\in\{0,1\}^n}
\mathsf T\!\left(
\bigl(o(x_i,b_i,y_i;u_i)\bigr)_{i=1}^n,\mathcal V
\right),
\]
where
\[
\mathbf z=((x_i,b_i,y_i))_{i=1}^n,
\qquad
\mathbf u=(u_i)_{i=1}^n,
\qquad
\mathcal V\in\mathscr B([-1,1]),
\]
and \(o(x,b,y;u)\) is the synthetic observation map of
\cref{obj:lem:synthetic-randomization-kernel}. This finite mixture composes the parameter-independent fair-coin transformation with \(\mathsf T\), and hence is an admissible observational kernel.

For every \(P_Z\in\mathcal D_{d,q}\),
\cref{obj:lem:synthetic-randomization-kernel} supplies a completion
\(P_Z^*\in\mathcal M^+_{n,d,q}\) whose observed product law is exactly the transformed sample law and whose target satisfies
\[
\tau(P_Z^*)=\theta_Z(P_Z).
\]
Thus
\[
\mathcal R^Z_{P_Z}(\widehat\theta_{\mathsf T})
=
\mathcal R_{P_Z^*}(\mathsf T).
\]
Both targets and all kernel outputs lie in \([-1,1]\), giving
\[
0\le\mathcal R^Z_{P_Z}(\widehat\theta_{\mathsf T})\le4,
\qquad
0\le\mathcal R_P(\mathsf T)\le4.
\]
The supremum comparison and the observational minimax definition therefore give, for every \(\mathsf T\),
\[
\mathfrak R^Z_{n,d,q}
\le
\sup_{P_Z\in\mathcal D_{d,q}}
\mathcal R^Z_{P_Z}(\widehat\theta_{\mathsf T})
\le
\sup_{P\in\mathcal M^+_{n,d,q}}
\mathcal R_P(\mathsf T).
\]
Taking the infimum over \(\mathsf T\) yields
\[
\mathfrak R^Z_{n,d,q}\le\mathfrak R^+_{n,d,q}.
\]

Now define
\[
\underline c(q):=\frac{\underline c_Z(q)}{\beta_q^2}>0.
\]
Using the rate comparison established in the preceding step,
\[
\underline c(q)r_{n,d,q}
\le
\frac{\underline c_Z(q)}{\beta_q^2}
\,\beta_q^2\mathfrak h_{n,d}
=
\underline c_Z(q)\mathfrak h_{n,d}
\le
\mathfrak R^Z_{n,d,q}
\le
\mathfrak R^+_{n,d,q}.
\]
% lean: CausalSmith.Stat.MarRareqLogfrontier.zeng_fixed_positivity_polynomial_frontier; CausalSmith.Stat.MarRareqLogfrontier.zengMinimaxRisk_le_unrestrictedMinimaxRisk; CausalSmith.Stat.MarRareqLogfrontier.unrestricted_fixed_q_minimax_frontier

\item
Choose the common upper constant
\[
\overline C(q)
:=
\max\{\overline C_{\mathrm{mix}},
       \overline C_{\mathrm{poly}}(q)\}
+\underline c(q).
\]
These are finite real constants depending only on \(q\), and
\[
0<\underline c(q)\le\overline C(q).
\]
The rates are nonnegative:
\[
\mathfrak h_{n,d}\ge0,
\qquad r_{n,d,q}\ge0,
\]
and the chosen constant satisfies
\[
\overline C_{\mathrm{mix}}\le\overline C(q),
\qquad
\overline C_{\mathrm{poly}}(q)\le\overline C(q).
\]

To obtain the fixed-rate lower bound with the same constant, observe that
\[
0<N\le n,
\qquad
0<\ell\le\ell_{n,1},
\qquad
0<N\ell\le n\ell_{n,1}.
\]
It follows that
\[
\frac1n\le\frac1N,
\qquad
0\le\frac{d}{n\ell_{n,1}}\le\frac{d}{N\ell},
\]
and hence
\[
\left(\frac{d}{n\ell_{n,1}}\right)^2
\le
\left(\frac{d}{N\ell}\right)^2.
\]
Adding and truncating at one proves
\[
\mathfrak h_{n,d}
=
\min\left\{1,\frac1n+
\left(\frac{d}{n\ell_{n,1}}\right)^2\right\}
\le
\min\left\{1,\frac1N+
\left(\frac{d}{N\ell}\right)^2\right\}
=r_{n,d,q}.
\]
Combining this comparison with the preceding bounds gives
\[
\begin{aligned}
\underline c(q)\mathfrak h_{n,d}
&\le \underline c(q)r_{n,d,q}
 \le\mathfrak R^+_{n,d,q},\\
\mathfrak R^+_{n,d,q}
&\le\overline C_{\mathrm{poly}}(q)\mathfrak h_{n,d}
 \le\overline C(q)\mathfrak h_{n,d},\\
\underline c(q)r_{n,d,q}
&\le\mathfrak R^+_{n,d,q},\\
\mathfrak R^+_{n,d,q}
&\le\overline C_{\mathrm{mix}}r_{n,d,q}
 \le\overline C(q)r_{n,d,q}.
\end{aligned}
\]
These are the two asserted comparisons.
% lean: CausalSmith.Stat.MarRareqLogfrontier.fixedQRate_le_frontierRate; CausalSmith.Stat.MarRareqLogfrontier.unrestricted_fixed_q_minimax_frontier
\end{enumerate}
\end{proof}
